%!TEX TS-program = lualatex
%!TEX encoding = UTF-8
% Gabarit d'un tableau périodique TikZ : coordonnées des cases, styles et légende.
% Les cases, symboles et informations viennent de simple_metaux.tex, que produit
% simple_metaux.py (pychimie.TableauTikz) :
%     python simple_metaux.py && lualatex Tableau_simple_metaux.tex
\documentclass{standalone}
\usepackage{tikz}
\usepackage{siunitx}
\usepackage{hyperref}

\hypersetup{
	colorlinks = false,
	urlbordercolor = .1 .1 1,
	pdfborderstyle = {/S/U/W 1} % liens soulignés, ne marche qu'avec acrobat
}

\newcommand\dessinecase[2][]{\draw[#1] (#2) ++(-.48, -.48) rectangle ++(.96, .96);}
\newcommand\baselegend[4]{\draw[#1] (Legend) ++ ({#2*\larglegend}, -#3*\hautlegend) coordinate(Ltmp) rectangle ++(.5, .25) #4; }
\newcommand\legend[4]{\baselegend{#1}{#2}{#3}{ (Ltmp) ++ (0,.125) node[legend, left]{#4} ; }}
\newcommand\legendhoriz[4]{\baselegend{#1}{#2}{#3}{ (Ltmp) ++ (.25,0) node[legend, below]{#4} ; }}

\def\larg{1cm}
\def\haut{1cm}
\def\xoriginlanth{3.5}
\def\yoriginlanth{-1}
\def\xoriginlegend{1.8}
\def\yoriginlegend{-3.6}
\def\larglegend{1.5}
\def\hautlegend{.4}

\sisetup{
	round-mode = places, % masses arrondies au centième
	round-precision = 2,
	output-decimal-marker = {,},
}

\begin{document}
\fontsize{9}{9}\selectfont

\begin{tikzpicture}[x=\larg, y=\haut,
	% une valeur absente ne produit pas de clé : seules les valeurs présentes
	% dans la table ont besoin d'un style
	etat/.is choice,
	etat/g/.style={text=red!80!white},
	etat/l/.style={text=blue!80!white},
	etat/s/.style={black},
	famille/.is choice,
	famille/métal alcalin/.style={fill=red!50!white},
	famille/métal alcalino-terreux/.style={fill=yellow!80!black},
	famille/lanthanide/.style={fill=red!20!white!95!blue},
	famille/actinide/.style={fill=red!30!white},
	famille/métal de transition/.style={fill=red!20!white},
	famille/métal pauvre/.style={fill=gray!60!white},
	famille/métalloïde/.style={fill=red!25!green!70!blue!25!gray},
	famille/non-métal/.style={fill=green!80!yellow!70!blue!60!white},
	famille/halogène/.style={fill=yellow!80!white!90!red},
	famille/gaz noble/.style={fill=blue!20!white!90!green},
	bloc/.is choice,
	bloc/s/.style={},
	bloc/p/.style={},
	bloc/d/.style={},
	bloc/f/.style={},
	titre colonne/.style={gray, scale=.7, node font={\itshape\footnotesize}, yshift=-.4\haut},
	titre ligne/.style={gray, scale=.7, node font={\itshape\footnotesize}, xshift=.4\larg},
	symbole/.style={node font={\bfseries\normalsize}, yshift=.05\haut},
	legend/.style={scale=.9, node font={\footnotesize}, text centered, text width=1.5cm},
	M info/.style={node font={\tiny}, yshift=-.3\haut},
	nom/.style={scale=.7, node font={\tiny}, yshift=.5\haut, text centered},
	Z/.style={node font={\bfseries\tiny}, yshift=.2\haut, xshift=.3\larg}
]
\coordinate (Legend) at (\xoriginlegend, \yoriginlegend);

% coordonnées <groupe>-<période> du tableau principal
\foreach \col in {1,...,18} {
	\foreach \ligne in {1,...,7} {
		\draw (\col, 7-\ligne) coordinate(\col-\ligne);
	}
}

% titres
\foreach \col in {1,18} \draw (\col, 7) node[titre colonne]{\col};
\foreach \col in {2, 13, 14, 15, 16, 17} \draw (\col, 6) node[titre colonne]{\col};
\foreach \col in {3, ..., 12} \draw (\col, 4) node[titre colonne]{\col};
\foreach \ligne in {1,...,7} \draw (0, {(7-\ligne)}) node[titre ligne]{\ligne};

% coordonnées <rang>-L et <rang>-A des lignes du bloc f
\foreach \col in {1,...,14} {
	\draw (\xoriginlanth-1, \yoriginlanth) ++ (\col, 0) node[titre colonne]{\col};
	\draw ({\xoriginlanth+\col-1}, {\yoriginlanth-1}) coordinate(\col-L);
	\draw ({\xoriginlanth+\col-1}, {\yoriginlanth-2}) coordinate(\col-A);
}

% éléments
\dessinecase[famille=non-métal]{1-1}
\dessinecase[famille=gaz noble]{18-1}
\dessinecase[famille=métal alcalin]{1-2}
\dessinecase[famille=métal alcalino-terreux]{2-2}
\dessinecase[famille=métalloïde]{13-2}
\dessinecase[famille=non-métal]{14-2}
\dessinecase[famille=non-métal]{15-2}
\dessinecase[famille=non-métal]{16-2}
\dessinecase[famille=halogène]{17-2}
\dessinecase[famille=gaz noble]{18-2}
\dessinecase[famille=métal alcalin]{1-3}
\dessinecase[famille=métal alcalino-terreux]{2-3}
\dessinecase[famille=métal pauvre]{13-3}
\dessinecase[famille=métalloïde]{14-3}
\dessinecase[famille=non-métal]{15-3}
\dessinecase[famille=non-métal]{16-3}
\dessinecase[famille=halogène]{17-3}
\dessinecase[famille=gaz noble]{18-3}
\dessinecase[famille=métal alcalin]{1-4}
\dessinecase[famille=métal alcalino-terreux]{2-4}
\dessinecase[famille=métal de transition]{3-4}
\dessinecase[famille=métal de transition]{4-4}
\dessinecase[famille=métal de transition]{5-4}
\dessinecase[famille=métal de transition]{6-4}
\dessinecase[famille=métal de transition]{7-4}
\dessinecase[famille=métal de transition]{8-4}
\dessinecase[famille=métal de transition]{9-4}
\dessinecase[famille=métal de transition]{10-4}
\dessinecase[famille=métal de transition]{11-4}
\dessinecase[famille=métal de transition]{12-4}
\dessinecase[famille=métal pauvre]{13-4}
\dessinecase[famille=métalloïde]{14-4}
\dessinecase[famille=métalloïde]{15-4}
\dessinecase[famille=non-métal]{16-4}
\dessinecase[famille=halogène]{17-4}
\dessinecase[famille=gaz noble]{18-4}
\dessinecase[famille=métal alcalin]{1-5}
\dessinecase[famille=métal alcalino-terreux]{2-5}
\dessinecase[famille=métal de transition]{3-5}
\dessinecase[famille=métal de transition]{4-5}
\dessinecase[famille=métal de transition]{5-5}
\dessinecase[famille=métal de transition]{6-5}
\dessinecase[famille=métal de transition]{7-5}
\dessinecase[famille=métal de transition]{8-5}
\dessinecase[famille=métal de transition]{9-5}
\dessinecase[famille=métal de transition]{10-5}
\dessinecase[famille=métal de transition]{11-5}
\dessinecase[famille=métal de transition]{12-5}
\dessinecase[famille=métal pauvre]{13-5}
\dessinecase[famille=métal pauvre]{14-5}
\dessinecase[famille=métalloïde]{15-5}
\dessinecase[famille=métalloïde]{16-5}
\dessinecase[famille=halogène]{17-5}
\dessinecase[famille=gaz noble]{18-5}
\dessinecase[famille=métal alcalin]{1-6}
\dessinecase[famille=métal alcalino-terreux]{2-6}
\dessinecase[famille=lanthanide]{1-L}
\dessinecase[famille=lanthanide]{2-L}
\dessinecase[famille=lanthanide]{3-L}
\dessinecase[famille=lanthanide]{4-L}
\dessinecase[famille=lanthanide]{5-L}
\dessinecase[famille=lanthanide]{6-L}
\dessinecase[famille=lanthanide]{7-L}
\dessinecase[famille=lanthanide]{8-L}
\dessinecase[famille=lanthanide]{9-L}
\dessinecase[famille=lanthanide]{10-L}
\dessinecase[famille=lanthanide]{11-L}
\dessinecase[famille=lanthanide]{12-L}
\dessinecase[famille=lanthanide]{13-L}
\dessinecase[famille=lanthanide]{14-L}
\dessinecase[famille=lanthanide]{3-6}
\dessinecase[famille=métal de transition]{4-6}
\dessinecase[famille=métal de transition]{5-6}
\dessinecase[famille=métal de transition]{6-6}
\dessinecase[famille=métal de transition]{7-6}
\dessinecase[famille=métal de transition]{8-6}
\dessinecase[famille=métal de transition]{9-6}
\dessinecase[famille=métal de transition]{10-6}
\dessinecase[famille=métal de transition]{11-6}
\dessinecase[famille=métal de transition]{12-6}
\dessinecase[famille=métal pauvre]{13-6}
\dessinecase[famille=métal pauvre]{14-6}
\dessinecase[famille=métal pauvre]{15-6}
\dessinecase[famille=métalloïde]{16-6}
\dessinecase[famille=halogène]{17-6}
\dessinecase[famille=gaz noble]{18-6}
\dessinecase[famille=métal alcalin]{1-7}
\dessinecase[famille=métal alcalino-terreux]{2-7}
\dessinecase[famille=actinide]{1-A}
\dessinecase[famille=actinide]{2-A}
\dessinecase[famille=actinide]{3-A}
\dessinecase[famille=actinide]{4-A}
\dessinecase[famille=actinide]{5-A}
\dessinecase[famille=actinide]{6-A}
\dessinecase[famille=actinide]{7-A}
\dessinecase[famille=actinide]{8-A}
\dessinecase[famille=actinide]{9-A}
\dessinecase[famille=actinide]{10-A}
\dessinecase[famille=actinide]{11-A}
\dessinecase[famille=actinide]{12-A}
\dessinecase[famille=actinide]{13-A}
\dessinecase[famille=actinide]{14-A}
\dessinecase[famille=actinide]{3-7}
\dessinecase[famille=métal de transition]{4-7}
\dessinecase[famille=métal de transition]{5-7}
\dessinecase[famille=métal de transition]{6-7}
\dessinecase[famille=métal de transition]{7-7}
\dessinecase[famille=métal de transition]{8-7}
\dessinecase[famille=métal de transition]{9-7}
\dessinecase[famille=métal de transition]{10-7}
\dessinecase[famille=métal de transition]{11-7}
\dessinecase[famille=métal de transition]{12-7}
\dessinecase[famille=métal pauvre]{13-7}
\dessinecase[famille=métal pauvre]{14-7}
\dessinecase[famille=métal pauvre]{15-7}
\dessinecase[famille=métal pauvre]{16-7}
\dessinecase[famille=halogène]{17-7}
\dessinecase[famille=gaz noble]{18-7}
\draw (1-1) coordinate(H) node[symbole, etat=g] {\href{https://fr.wikipedia.org/wiki/Hydrog\%C3\%A8ne}{H}};
\draw (18-1) coordinate(He) node[symbole, etat=g] {\href{https://fr.wikipedia.org/wiki/H\%C3\%A9lium}{He}};
\draw (1-2) coordinate(Li) node[symbole, etat=s] {\href{https://fr.wikipedia.org/wiki/Lithium}{Li}};
\draw (2-2) coordinate(Be) node[symbole, etat=s] {\href{https://fr.wikipedia.org/wiki/B\%C3\%A9ryllium}{Be}};
\draw (13-2) coordinate(B) node[symbole, etat=s] {\href{https://fr.wikipedia.org/wiki/Bore}{B}};
\draw (14-2) coordinate(C) node[symbole, etat=s] {\href{https://fr.wikipedia.org/wiki/Carbone}{C}};
\draw (15-2) coordinate(N) node[symbole, etat=g] {\href{https://fr.wikipedia.org/wiki/Azote}{N}};
\draw (16-2) coordinate(O) node[symbole, etat=g] {\href{https://fr.wikipedia.org/wiki/Oxyg\%C3\%A8ne}{O}};
\draw (17-2) coordinate(F) node[symbole, etat=g] {\href{https://fr.wikipedia.org/wiki/Fluor}{F}};
\draw (18-2) coordinate(Ne) node[symbole, etat=g] {\href{https://fr.wikipedia.org/wiki/N\%C3\%A9on}{Ne}};
\draw (1-3) coordinate(Na) node[symbole, etat=s] {\href{https://fr.wikipedia.org/wiki/Sodium}{Na}};
\draw (2-3) coordinate(Mg) node[symbole, etat=s] {\href{https://fr.wikipedia.org/wiki/Magn\%C3\%A9sium}{Mg}};
\draw (13-3) coordinate(Al) node[symbole, etat=s] {\href{https://fr.wikipedia.org/wiki/Aluminium}{Al}};
\draw (14-3) coordinate(Si) node[symbole, etat=s] {\href{https://fr.wikipedia.org/wiki/Silicium}{Si}};
\draw (15-3) coordinate(P) node[symbole, etat=s] {\href{https://fr.wikipedia.org/wiki/Phosphore}{P}};
\draw (16-3) coordinate(S) node[symbole, etat=s] {\href{https://fr.wikipedia.org/wiki/Soufre}{S}};
\draw (17-3) coordinate(Cl) node[symbole, etat=g] {\href{https://fr.wikipedia.org/wiki/Chlore}{Cl}};
\draw (18-3) coordinate(Ar) node[symbole, etat=g] {\href{https://fr.wikipedia.org/wiki/Argon}{Ar}};
\draw (1-4) coordinate(K) node[symbole, etat=s] {\href{https://fr.wikipedia.org/wiki/Potassium}{K}};
\draw (2-4) coordinate(Ca) node[symbole, etat=s] {\href{https://fr.wikipedia.org/wiki/Calcium}{Ca}};
\draw (3-4) coordinate(Sc) node[symbole, etat=s] {\href{https://fr.wikipedia.org/wiki/Scandium}{Sc}};
\draw (4-4) coordinate(Ti) node[symbole, etat=s] {\href{https://fr.wikipedia.org/wiki/Titane}{Ti}};
\draw (5-4) coordinate(V) node[symbole, etat=s] {\href{https://fr.wikipedia.org/wiki/Vanadium}{V}};
\draw (6-4) coordinate(Cr) node[symbole, etat=s] {\href{https://fr.wikipedia.org/wiki/Chrome}{Cr}};
\draw (7-4) coordinate(Mn) node[symbole, etat=s] {\href{https://fr.wikipedia.org/wiki/Mangan\%C3\%A8se}{Mn}};
\draw (8-4) coordinate(Fe) node[symbole, etat=s] {\href{https://fr.wikipedia.org/wiki/Fer}{Fe}};
\draw (9-4) coordinate(Co) node[symbole, etat=s] {\href{https://fr.wikipedia.org/wiki/Cobalt}{Co}};
\draw (10-4) coordinate(Ni) node[symbole, etat=s] {\href{https://fr.wikipedia.org/wiki/Nickel}{Ni}};
\draw (11-4) coordinate(Cu) node[symbole, etat=s] {\href{https://fr.wikipedia.org/wiki/Cuivre}{Cu}};
\draw (12-4) coordinate(Zn) node[symbole, etat=s] {\href{https://fr.wikipedia.org/wiki/Zinc}{Zn}};
\draw (13-4) coordinate(Ga) node[symbole, etat=s] {\href{https://fr.wikipedia.org/wiki/Gallium}{Ga}};
\draw (14-4) coordinate(Ge) node[symbole, etat=s] {\href{https://fr.wikipedia.org/wiki/Germanium}{Ge}};
\draw (15-4) coordinate(As) node[symbole, etat=s] {\href{https://fr.wikipedia.org/wiki/Arsenic}{As}};
\draw (16-4) coordinate(Se) node[symbole, etat=s] {\href{https://fr.wikipedia.org/wiki/S\%C3\%A9l\%C3\%A9nium}{Se}};
\draw (17-4) coordinate(Br) node[symbole, etat=l] {\href{https://fr.wikipedia.org/wiki/Brome}{Br}};
\draw (18-4) coordinate(Kr) node[symbole, etat=g] {\href{https://fr.wikipedia.org/wiki/Krypton}{Kr}};
\draw (1-5) coordinate(Rb) node[symbole, etat=s] {\href{https://fr.wikipedia.org/wiki/Rubidium}{Rb}};
\draw (2-5) coordinate(Sr) node[symbole, etat=s] {\href{https://fr.wikipedia.org/wiki/Strontium}{Sr}};
\draw (3-5) coordinate(Y) node[symbole, etat=s] {\href{https://fr.wikipedia.org/wiki/Yttrium}{Y}};
\draw (4-5) coordinate(Zr) node[symbole, etat=s] {\href{https://fr.wikipedia.org/wiki/Zirconium}{Zr}};
\draw (5-5) coordinate(Nb) node[symbole, etat=s] {\href{https://fr.wikipedia.org/wiki/Niobium}{Nb}};
\draw (6-5) coordinate(Mo) node[symbole, etat=s] {\href{https://fr.wikipedia.org/wiki/Molybd\%C3\%A8ne}{Mo}};
\draw (7-5) coordinate(Tc) node[symbole, etat=s] {\href{https://fr.wikipedia.org/wiki/Techn\%C3\%A9tium}{Tc}};
\draw (8-5) coordinate(Ru) node[symbole, etat=s] {\href{https://fr.wikipedia.org/wiki/Ruth\%C3\%A9nium}{Ru}};
\draw (9-5) coordinate(Rh) node[symbole, etat=s] {\href{https://fr.wikipedia.org/wiki/Rhodium}{Rh}};
\draw (10-5) coordinate(Pd) node[symbole, etat=s] {\href{https://fr.wikipedia.org/wiki/Palladium}{Pd}};
\draw (11-5) coordinate(Ag) node[symbole, etat=s] {\href{https://fr.wikipedia.org/wiki/Argent}{Ag}};
\draw (12-5) coordinate(Cd) node[symbole, etat=s] {\href{https://fr.wikipedia.org/wiki/Cadmium}{Cd}};
\draw (13-5) coordinate(In) node[symbole, etat=s] {\href{https://fr.wikipedia.org/wiki/Indium}{In}};
\draw (14-5) coordinate(Sn) node[symbole, etat=s] {\href{https://fr.wikipedia.org/wiki/\%C3\%89tain}{Sn}};
\draw (15-5) coordinate(Sb) node[symbole, etat=s] {\href{https://fr.wikipedia.org/wiki/Antimoine}{Sb}};
\draw (16-5) coordinate(Te) node[symbole, etat=s] {\href{https://fr.wikipedia.org/wiki/Tellure}{Te}};
\draw (17-5) coordinate(I) node[symbole, etat=s] {\href{https://fr.wikipedia.org/wiki/Iode}{I}};
\draw (18-5) coordinate(Xe) node[symbole, etat=g] {\href{https://fr.wikipedia.org/wiki/X\%C3\%A9non}{Xe}};
\draw (1-6) coordinate(Cs) node[symbole, etat=s] {\href{https://fr.wikipedia.org/wiki/C\%C3\%A9sium}{Cs}};
\draw (2-6) coordinate(Ba) node[symbole, etat=s] {\href{https://fr.wikipedia.org/wiki/Baryum}{Ba}};
\draw (1-L) coordinate(La) node[symbole, etat=s] {\href{https://fr.wikipedia.org/wiki/Lanthane}{La}};
\draw (2-L) coordinate(Ce) node[symbole, etat=s] {\href{https://fr.wikipedia.org/wiki/C\%C3\%A9rium}{Ce}};
\draw (3-L) coordinate(Pr) node[symbole, etat=s] {\href{https://fr.wikipedia.org/wiki/Pras\%C3\%A9odyme}{Pr}};
\draw (4-L) coordinate(Nd) node[symbole, etat=s] {\href{https://fr.wikipedia.org/wiki/N\%C3\%A9odyme}{Nd}};
\draw (5-L) coordinate(Pm) node[symbole, etat=s] {\href{https://fr.wikipedia.org/wiki/Prom\%C3\%A9thium}{Pm}};
\draw (6-L) coordinate(Sm) node[symbole, etat=s] {\href{https://fr.wikipedia.org/wiki/Samarium}{Sm}};
\draw (7-L) coordinate(Eu) node[symbole, etat=s] {\href{https://fr.wikipedia.org/wiki/Europium}{Eu}};
\draw (8-L) coordinate(Gd) node[symbole, etat=s] {\href{https://fr.wikipedia.org/wiki/Gadolinium}{Gd}};
\draw (9-L) coordinate(Tb) node[symbole, etat=s] {\href{https://fr.wikipedia.org/wiki/Terbium}{Tb}};
\draw (10-L) coordinate(Dy) node[symbole, etat=s] {\href{https://fr.wikipedia.org/wiki/Dysprosium}{Dy}};
\draw (11-L) coordinate(Ho) node[symbole, etat=s] {\href{https://fr.wikipedia.org/wiki/Holmium}{Ho}};
\draw (12-L) coordinate(Er) node[symbole, etat=s] {\href{https://fr.wikipedia.org/wiki/Erbium}{Er}};
\draw (13-L) coordinate(Tm) node[symbole, etat=s] {\href{https://fr.wikipedia.org/wiki/Thulium}{Tm}};
\draw (14-L) coordinate(Yb) node[symbole, etat=s] {\href{https://fr.wikipedia.org/wiki/Ytterbium}{Yb}};
\draw (3-6) coordinate(Lu) node[symbole, etat=s] {\href{https://fr.wikipedia.org/wiki/Lut\%C3\%A9cium}{Lu}};
\draw (4-6) coordinate(Hf) node[symbole, etat=s] {\href{https://fr.wikipedia.org/wiki/Hafnium}{Hf}};
\draw (5-6) coordinate(Ta) node[symbole, etat=s] {\href{https://fr.wikipedia.org/wiki/Tantale_\%28chimie\%29}{Ta}};
\draw (6-6) coordinate(W) node[symbole, etat=s] {\href{https://fr.wikipedia.org/wiki/Tungst\%C3\%A8ne}{W}};
\draw (7-6) coordinate(Re) node[symbole, etat=s] {\href{https://fr.wikipedia.org/wiki/Rh\%C3\%A9nium}{Re}};
\draw (8-6) coordinate(Os) node[symbole, etat=s] {\href{https://fr.wikipedia.org/wiki/Osmium}{Os}};
\draw (9-6) coordinate(Ir) node[symbole, etat=s] {\href{https://fr.wikipedia.org/wiki/Iridium}{Ir}};
\draw (10-6) coordinate(Pt) node[symbole, etat=s] {\href{https://fr.wikipedia.org/wiki/Platine}{Pt}};
\draw (11-6) coordinate(Au) node[symbole, etat=s] {\href{https://fr.wikipedia.org/wiki/Or}{Au}};
\draw (12-6) coordinate(Hg) node[symbole, etat=l] {\href{https://fr.wikipedia.org/wiki/Mercure_\%28chimie\%29}{Hg}};
\draw (13-6) coordinate(Tl) node[symbole, etat=s] {\href{https://fr.wikipedia.org/wiki/Thallium}{Tl}};
\draw (14-6) coordinate(Pb) node[symbole, etat=s] {\href{https://fr.wikipedia.org/wiki/Plomb}{Pb}};
\draw (15-6) coordinate(Bi) node[symbole, etat=s] {\href{https://fr.wikipedia.org/wiki/Bismuth}{Bi}};
\draw (16-6) coordinate(Po) node[symbole, etat=s] {\href{https://fr.wikipedia.org/wiki/Polonium}{Po}};
\draw (17-6) coordinate(At) node[symbole] {\href{https://fr.wikipedia.org/wiki/Astate}{At}};
\draw (18-6) coordinate(Rn) node[symbole, etat=g] {\href{https://fr.wikipedia.org/wiki/Radon}{Rn}};
\draw (1-7) coordinate(Fr) node[symbole] {\href{https://fr.wikipedia.org/wiki/Francium}{Fr}};
\draw (2-7) coordinate(Ra) node[symbole, etat=s] {\href{https://fr.wikipedia.org/wiki/Radium}{Ra}};
\draw (1-A) coordinate(Ac) node[symbole, etat=s] {\href{https://fr.wikipedia.org/wiki/Actinium}{Ac}};
\draw (2-A) coordinate(Th) node[symbole, etat=s] {\href{https://fr.wikipedia.org/wiki/Thorium}{Th}};
\draw (3-A) coordinate(Pa) node[symbole, etat=s] {\href{https://fr.wikipedia.org/wiki/Protactinium}{Pa}};
\draw (4-A) coordinate(U) node[symbole, etat=s] {\href{https://fr.wikipedia.org/wiki/Uranium}{U}};
\draw (5-A) coordinate(Np) node[symbole, etat=s] {\href{https://fr.wikipedia.org/wiki/Neptunium}{Np}};
\draw (6-A) coordinate(Pu) node[symbole, etat=s] {\href{https://fr.wikipedia.org/wiki/Plutonium}{Pu}};
\draw (7-A) coordinate(Am) node[symbole, etat=s] {\href{https://fr.wikipedia.org/wiki/Am\%C3\%A9ricium}{Am}};
\draw (8-A) coordinate(Cm) node[symbole, etat=s] {\href{https://fr.wikipedia.org/wiki/Curium}{Cm}};
\draw (9-A) coordinate(Bk) node[symbole, etat=s] {\href{https://fr.wikipedia.org/wiki/Berk\%C3\%A9lium}{Bk}};
\draw (10-A) coordinate(Cf) node[symbole, etat=s] {\href{https://fr.wikipedia.org/wiki/Californium}{Cf}};
\draw (11-A) coordinate(Es) node[symbole, etat=s] {\href{https://fr.wikipedia.org/wiki/Einsteinium}{Es}};
\draw (12-A) coordinate(Fm) node[symbole] {\href{https://fr.wikipedia.org/wiki/Fermium}{Fm}};
\draw (13-A) coordinate(Md) node[symbole] {\href{https://fr.wikipedia.org/wiki/Mend\%C3\%A9l\%C3\%A9vium}{Md}};
\draw (14-A) coordinate(No) node[symbole] {\href{https://fr.wikipedia.org/wiki/Nob\%C3\%A9lium}{No}};
\draw (3-7) coordinate(Lr) node[symbole] {\href{https://fr.wikipedia.org/wiki/Lawrencium}{Lr}};
\draw (4-7) coordinate(Rf) node[symbole] {\href{https://fr.wikipedia.org/wiki/Rutherfordium}{Rf}};
\draw (5-7) coordinate(Db) node[symbole] {\href{https://fr.wikipedia.org/wiki/Dubnium}{Db}};
\draw (6-7) coordinate(Sg) node[symbole] {\href{https://fr.wikipedia.org/wiki/Seaborgium}{Sg}};
\draw (7-7) coordinate(Bh) node[symbole] {\href{https://fr.wikipedia.org/wiki/Bohrium}{Bh}};
\draw (8-7) coordinate(Hs) node[symbole] {\href{https://fr.wikipedia.org/wiki/Hassium}{Hs}};
\draw (9-7) coordinate(Mt) node[symbole] {\href{https://fr.wikipedia.org/wiki/Meitn\%C3\%A9rium}{Mt}};
\draw (10-7) coordinate(Ds) node[symbole] {\href{https://fr.wikipedia.org/wiki/Darmstadtium}{Ds}};
\draw (11-7) coordinate(Rg) node[symbole] {\href{https://fr.wikipedia.org/wiki/Roentgenium}{Rg}};
\draw (12-7) coordinate(Cn) node[symbole] {\href{https://fr.wikipedia.org/wiki/Copernicium}{Cn}};
\draw (13-7) coordinate(Nh) node[symbole] {\href{https://fr.wikipedia.org/wiki/Nihonium}{Nh}};
\draw (14-7) coordinate(Fl) node[symbole] {\href{https://fr.wikipedia.org/wiki/Fl\%C3\%A9rovium}{Fl}};
\draw (15-7) coordinate(Mc) node[symbole] {\href{https://fr.wikipedia.org/wiki/Moscovium}{Mc}};
\draw (16-7) coordinate(Lv) node[symbole] {\href{https://fr.wikipedia.org/wiki/Livermorium}{Lv}};
\draw (17-7) coordinate(Ts) node[symbole] {\href{https://fr.wikipedia.org/wiki/Tennesse}{Ts}};
\draw (18-7) coordinate(Og) node[symbole] {\href{https://fr.wikipedia.org/wiki/Oganesson}{Og}};
\draw (1-1) node[M info] {\num{1.008}};
\draw (18-1) node[M info] {\num{4.0026}};
\draw (1-2) node[M info] {\num{6.94}};
\draw (2-2) node[M info] {\num{9.0122}};
\draw (13-2) node[M info] {\num{10.81}};
\draw (14-2) node[M info] {\num{12.011}};
\draw (15-2) node[M info] {\num{14.007}};
\draw (16-2) node[M info] {\num{15.999}};
\draw (17-2) node[M info] {\num{18.998}};
\draw (18-2) node[M info] {\num{20.18}};
\draw (1-3) node[M info] {\num{22.99}};
\draw (2-3) node[M info] {\num{24.305}};
\draw (13-3) node[M info] {\num{26.982}};
\draw (14-3) node[M info] {\num{28.085}};
\draw (15-3) node[M info] {\num{30.974}};
\draw (16-3) node[M info] {\num{32.06}};
\draw (17-3) node[M info] {\num{35.45}};
\draw (18-3) node[M info] {\num{39.95}};
\draw (1-4) node[M info] {\num{39.098}};
\draw (2-4) node[M info] {\num{40.078}};
\draw (3-4) node[M info] {\num{44.956}};
\draw (4-4) node[M info] {\num{47.867}};
\draw (5-4) node[M info] {\num{50.942}};
\draw (6-4) node[M info] {\num{51.996}};
\draw (7-4) node[M info] {\num{54.938}};
\draw (8-4) node[M info] {\num{55.845}};
\draw (9-4) node[M info] {\num{58.933}};
\draw (10-4) node[M info] {\num{58.693}};
\draw (11-4) node[M info] {\num{63.546}};
\draw (12-4) node[M info] {\num{65.38}};
\draw (13-4) node[M info] {\num{69.723}};
\draw (14-4) node[M info] {\num{72.63}};
\draw (15-4) node[M info] {\num{74.922}};
\draw (16-4) node[M info] {\num{78.971}};
\draw (17-4) node[M info] {\num{79.904}};
\draw (18-4) node[M info] {\num{83.798}};
\draw (1-5) node[M info] {\num{85.468}};
\draw (2-5) node[M info] {\num{87.62}};
\draw (3-5) node[M info] {\num{88.906}};
\draw (4-5) node[M info] {\num{91.222}};
\draw (5-5) node[M info] {\num{92.906}};
\draw (6-5) node[M info] {\num{95.95}};
\draw (8-5) node[M info] {\num{101.07}};
\draw (9-5) node[M info] {\num{102.91}};
\draw (10-5) node[M info] {\num{106.42}};
\draw (11-5) node[M info] {\num{107.87}};
\draw (12-5) node[M info] {\num{112.41}};
\draw (13-5) node[M info] {\num{114.82}};
\draw (14-5) node[M info] {\num{118.71}};
\draw (15-5) node[M info] {\num{121.76}};
\draw (16-5) node[M info] {\num{127.6}};
\draw (17-5) node[M info] {\num{126.9}};
\draw (18-5) node[M info] {\num{131.29}};
\draw (1-6) node[M info] {\num{132.91}};
\draw (2-6) node[M info] {\num{137.33}};
\draw (1-L) node[M info] {\num{138.91}};
\draw (2-L) node[M info] {\num{140.12}};
\draw (3-L) node[M info] {\num{140.91}};
\draw (4-L) node[M info] {\num{144.24}};
\draw (6-L) node[M info] {\num{150.36}};
\draw (7-L) node[M info] {\num{151.96}};
\draw (8-L) node[M info] {\num{157.25}};
\draw (9-L) node[M info] {\num{158.93}};
\draw (10-L) node[M info] {\num{162.5}};
\draw (11-L) node[M info] {\num{164.93}};
\draw (12-L) node[M info] {\num{167.26}};
\draw (13-L) node[M info] {\num{168.93}};
\draw (14-L) node[M info] {\num{173.05}};
\draw (3-6) node[M info] {\num{174.97}};
\draw (4-6) node[M info] {\num{178.49}};
\draw (5-6) node[M info] {\num{180.95}};
\draw (6-6) node[M info] {\num{183.84}};
\draw (7-6) node[M info] {\num{186.21}};
\draw (8-6) node[M info] {\num{190.23}};
\draw (9-6) node[M info] {\num{192.22}};
\draw (10-6) node[M info] {\num{195.08}};
\draw (11-6) node[M info] {\num{196.97}};
\draw (12-6) node[M info] {\num{200.59}};
\draw (13-6) node[M info] {\num{204.38}};
\draw (14-6) node[M info] {\num{207.2}};
\draw (15-6) node[M info] {\num{208.98}};
\draw (2-A) node[M info] {\num{232.04}};
\draw (3-A) node[M info] {\num{231.04}};
\draw (4-A) node[M info] {\num{238.03}};
\draw (1-1) node[nom] {Hydrogène};
\draw (18-1) node[nom] {Hélium};
\draw (1-2) node[nom] {Lithium};
\draw (2-2) node[nom] {Béryllium};
\draw (13-2) node[nom] {Bore};
\draw (14-2) node[nom] {Carbone};
\draw (15-2) node[nom] {Azote};
\draw (16-2) node[nom] {Oxygène};
\draw (17-2) node[nom] {Fluor};
\draw (18-2) node[nom] {Néon};
\draw (1-3) node[nom] {Sodium};
\draw (2-3) node[nom] {Magnésium};
\draw (13-3) node[nom] {Aluminium};
\draw (14-3) node[nom] {Silicium};
\draw (15-3) node[nom] {Phosphore};
\draw (16-3) node[nom] {Soufre};
\draw (17-3) node[nom] {Chlore};
\draw (18-3) node[nom] {Argon};
\draw (1-4) node[nom] {Potassium};
\draw (2-4) node[nom] {Calcium};
\draw (3-4) node[nom] {Scandium};
\draw (4-4) node[nom] {Titane};
\draw (5-4) node[nom] {Vanadium};
\draw (6-4) node[nom] {Chrome};
\draw (7-4) node[nom] {Manganèse};
\draw (8-4) node[nom] {Fer};
\draw (9-4) node[nom] {Cobalt};
\draw (10-4) node[nom] {Nickel};
\draw (11-4) node[nom] {Cuivre};
\draw (12-4) node[nom] {Zinc};
\draw (13-4) node[nom] {Gallium};
\draw (14-4) node[nom] {Germanium};
\draw (15-4) node[nom] {Arsenic};
\draw (16-4) node[nom] {Sélénium};
\draw (17-4) node[nom] {Brome};
\draw (18-4) node[nom] {Krypton};
\draw (1-5) node[nom] {Rubidium};
\draw (2-5) node[nom] {Strontium};
\draw (3-5) node[nom] {Yttrium};
\draw (4-5) node[nom] {Zirconium};
\draw (5-5) node[nom] {Niobium};
\draw (6-5) node[nom] {Molybdène};
\draw (7-5) node[nom] {Technétium};
\draw (8-5) node[nom] {Ruthénium};
\draw (9-5) node[nom] {Rhodium};
\draw (10-5) node[nom] {Palladium};
\draw (11-5) node[nom] {Argent};
\draw (12-5) node[nom] {Cadmium};
\draw (13-5) node[nom] {Indium};
\draw (14-5) node[nom] {Étain};
\draw (15-5) node[nom] {Antimoine};
\draw (16-5) node[nom] {Tellure};
\draw (17-5) node[nom] {Iode};
\draw (18-5) node[nom] {Xénon};
\draw (1-6) node[nom] {Césium};
\draw (2-6) node[nom] {Baryum};
\draw (1-L) node[nom] {Lanthane};
\draw (2-L) node[nom] {Cérium};
\draw (3-L) node[nom] {Praséodyme};
\draw (4-L) node[nom] {Néodyme};
\draw (5-L) node[nom] {Prométhium};
\draw (6-L) node[nom] {Samarium};
\draw (7-L) node[nom] {Europium};
\draw (8-L) node[nom] {Gadolinium};
\draw (9-L) node[nom] {Terbium};
\draw (10-L) node[nom] {Dysprosium};
\draw (11-L) node[nom] {Holmium};
\draw (12-L) node[nom] {Erbium};
\draw (13-L) node[nom] {Thulium};
\draw (14-L) node[nom] {Ytterbium};
\draw (3-6) node[nom] {Lutécium};
\draw (4-6) node[nom] {Hafnium};
\draw (5-6) node[nom] {Tantale};
\draw (6-6) node[nom] {Tungstène};
\draw (7-6) node[nom] {Rhénium};
\draw (8-6) node[nom] {Osmium};
\draw (9-6) node[nom] {Iridium};
\draw (10-6) node[nom] {Platine};
\draw (11-6) node[nom] {Or};
\draw (12-6) node[nom] {Mercure};
\draw (13-6) node[nom] {Thallium};
\draw (14-6) node[nom] {Plomb};
\draw (15-6) node[nom] {Bismuth};
\draw (16-6) node[nom] {Polonium};
\draw (17-6) node[nom] {Astate};
\draw (18-6) node[nom] {Radon};
\draw (1-7) node[nom] {Francium};
\draw (2-7) node[nom] {Radium};
\draw (1-A) node[nom] {Actinium};
\draw (2-A) node[nom] {Thorium};
\draw (3-A) node[nom] {Protactinium};
\draw (4-A) node[nom] {Uranium};
\draw (5-A) node[nom] {Neptunium};
\draw (6-A) node[nom] {Plutonium};
\draw (7-A) node[nom] {Américium};
\draw (8-A) node[nom] {Curium};
\draw (9-A) node[nom] {Berkélium};
\draw (10-A) node[nom] {Californium};
\draw (11-A) node[nom] {Einsteinium};
\draw (12-A) node[nom] {Fermium};
\draw (13-A) node[nom] {Mendélévium};
\draw (14-A) node[nom] {Nobélium};
\draw (3-7) node[nom] {Lawrencium};
\draw (4-7) node[nom] {Rutherfordium};
\draw (5-7) node[nom] {Dubnium};
\draw (6-7) node[nom] {Seaborgium};
\draw (7-7) node[nom] {Bohrium};
\draw (8-7) node[nom] {Hassium};
\draw (9-7) node[nom] {Meitnérium};
\draw (10-7) node[nom] {Darmstadtium};
\draw (11-7) node[nom] {Roentgenium};
\draw (12-7) node[nom] {Copernicium};
\draw (13-7) node[nom] {Nihonium};
\draw (14-7) node[nom] {Flérovium};
\draw (15-7) node[nom] {Moscovium};
\draw (16-7) node[nom] {Livermorium};
\draw (17-7) node[nom] {Tennesse};
\draw (18-7) node[nom] {Oganesson};
\draw (1-1) node[Z] {1};
\draw (18-1) node[Z] {2};
\draw (1-2) node[Z] {3};
\draw (2-2) node[Z] {4};
\draw (13-2) node[Z] {5};
\draw (14-2) node[Z] {6};
\draw (15-2) node[Z] {7};
\draw (16-2) node[Z] {8};
\draw (17-2) node[Z] {9};
\draw (18-2) node[Z] {10};
\draw (1-3) node[Z] {11};
\draw (2-3) node[Z] {12};
\draw (13-3) node[Z] {13};
\draw (14-3) node[Z] {14};
\draw (15-3) node[Z] {15};
\draw (16-3) node[Z] {16};
\draw (17-3) node[Z] {17};
\draw (18-3) node[Z] {18};
\draw (1-4) node[Z] {19};
\draw (2-4) node[Z] {20};
\draw (3-4) node[Z] {21};
\draw (4-4) node[Z] {22};
\draw (5-4) node[Z] {23};
\draw (6-4) node[Z] {24};
\draw (7-4) node[Z] {25};
\draw (8-4) node[Z] {26};
\draw (9-4) node[Z] {27};
\draw (10-4) node[Z] {28};
\draw (11-4) node[Z] {29};
\draw (12-4) node[Z] {30};
\draw (13-4) node[Z] {31};
\draw (14-4) node[Z] {32};
\draw (15-4) node[Z] {33};
\draw (16-4) node[Z] {34};
\draw (17-4) node[Z] {35};
\draw (18-4) node[Z] {36};
\draw (1-5) node[Z] {37};
\draw (2-5) node[Z] {38};
\draw (3-5) node[Z] {39};
\draw (4-5) node[Z] {40};
\draw (5-5) node[Z] {41};
\draw (6-5) node[Z] {42};
\draw (7-5) node[Z] {43};
\draw (8-5) node[Z] {44};
\draw (9-5) node[Z] {45};
\draw (10-5) node[Z] {46};
\draw (11-5) node[Z] {47};
\draw (12-5) node[Z] {48};
\draw (13-5) node[Z] {49};
\draw (14-5) node[Z] {50};
\draw (15-5) node[Z] {51};
\draw (16-5) node[Z] {52};
\draw (17-5) node[Z] {53};
\draw (18-5) node[Z] {54};
\draw (1-6) node[Z] {55};
\draw (2-6) node[Z] {56};
\draw (1-L) node[Z] {57};
\draw (2-L) node[Z] {58};
\draw (3-L) node[Z] {59};
\draw (4-L) node[Z] {60};
\draw (5-L) node[Z] {61};
\draw (6-L) node[Z] {62};
\draw (7-L) node[Z] {63};
\draw (8-L) node[Z] {64};
\draw (9-L) node[Z] {65};
\draw (10-L) node[Z] {66};
\draw (11-L) node[Z] {67};
\draw (12-L) node[Z] {68};
\draw (13-L) node[Z] {69};
\draw (14-L) node[Z] {70};
\draw (3-6) node[Z] {71};
\draw (4-6) node[Z] {72};
\draw (5-6) node[Z] {73};
\draw (6-6) node[Z] {74};
\draw (7-6) node[Z] {75};
\draw (8-6) node[Z] {76};
\draw (9-6) node[Z] {77};
\draw (10-6) node[Z] {78};
\draw (11-6) node[Z] {79};
\draw (12-6) node[Z] {80};
\draw (13-6) node[Z] {81};
\draw (14-6) node[Z] {82};
\draw (15-6) node[Z] {83};
\draw (16-6) node[Z] {84};
\draw (17-6) node[Z] {85};
\draw (18-6) node[Z] {86};
\draw (1-7) node[Z] {87};
\draw (2-7) node[Z] {88};
\draw (1-A) node[Z] {89};
\draw (2-A) node[Z] {90};
\draw (3-A) node[Z] {91};
\draw (4-A) node[Z] {92};
\draw (5-A) node[Z] {93};
\draw (6-A) node[Z] {94};
\draw (7-A) node[Z] {95};
\draw (8-A) node[Z] {96};
\draw (9-A) node[Z] {97};
\draw (10-A) node[Z] {98};
\draw (11-A) node[Z] {99};
\draw (12-A) node[Z] {100};
\draw (13-A) node[Z] {101};
\draw (14-A) node[Z] {102};
\draw (3-7) node[Z] {103};
\draw (4-7) node[Z] {104};
\draw (5-7) node[Z] {105};
\draw (6-7) node[Z] {106};
\draw (7-7) node[Z] {107};
\draw (8-7) node[Z] {108};
\draw (9-7) node[Z] {109};
\draw (10-7) node[Z] {110};
\draw (11-7) node[Z] {111};
\draw (12-7) node[Z] {112};
\draw (13-7) node[Z] {113};
\draw (14-7) node[Z] {114};
\draw (15-7) node[Z] {115};
\draw (16-7) node[Z] {116};
\draw (17-7) node[Z] {117};
\draw (18-7) node[Z] {118};
\legendhoriz{famille=métal alcalin}{0}{1}{Métal alcalin}
\legendhoriz{famille=métal alcalino-terreux}{1}{1}{Métal alcalino-terreux}
\legendhoriz{famille=lanthanide}{2}{1}{Lanthanide}
\legendhoriz{famille=actinide}{3}{1}{Actinide}
\legendhoriz{famille=métal de transition}{4}{1}{Métal de transition}
\legendhoriz{famille=métal pauvre}{5}{1}{Métal pauvre}
\legendhoriz{famille=métalloïde}{6}{1}{Métalloïde}
\legendhoriz{famille=non-métal}{7}{1}{Non-métal}
\legendhoriz{famille=halogène}{8}{1}{Halogène}
\legendhoriz{famille=gaz noble}{9}{1}{Gaz noble}


% renvoi des lignes du bloc f, entre les groupes 2 et 3
\draw[very thick] (2.5, -.5) -- ++(0,2) (\xoriginlanth-.5, \yoriginlanth-.5) -- ++(0,-2);
\draw (2.5, -.5) |- (\xoriginlanth-.5, \yoriginlanth-1.5);
\end{tikzpicture}

\end{document}
